\documentclass[aspectratio=169]{beamer}

% Tema y configuración
\usetheme{Madrid}
\usecolortheme{default}

% Paquetes necesarios
\usepackage[utf8]{inputenc}
\usepackage[spanish]{babel}
\usepackage{amsmath}
\usepackage{amsfonts}
\usepackage{amssymb}
\usepackage{graphicx}
\usepackage{listings}
% Habilitar UTF-8 dentro de lstlisting (acentos, ñ, etc.)
\usepackage{listingsutf8}
\usepackage{xcolor}
\usepackage{verbatim}
% Para dibujos con TikZ
\usepackage{tikz}
\usetikzlibrary{positioning,trees}

% Configuración de listings para código
\lstset{
    basicstyle=\ttfamily\footnotesize,
    keywordstyle=\color{blue}\bfseries,
    commentstyle=\color{green},
    stringstyle=\color{red},
    showstringspaces=false,
    breaklines=true,
    frame=single,
    numbers=left,
    numberstyle=\tiny\color{gray},
    inputencoding=utf8,
    % Mapeo de caracteres españoles para que no falle la compilación
    literate={á}{{\'a}}1 {é}{{\'e}}1 {í}{{\'\i}}1 {ó}{{\'o}}1 {ú}{{\'u}}1
             {Á}{{\'A}}1 {É}{{\'E}}1 {Í}{{\'I}}1 {Ó}{{\'O}}1 {Ú}{{\'U}}1
             {ñ}{{\~n}}1 {Ñ}{{\~N}}1 {ü}{{\"u}}1 {Ü}{{\"U}}1
}

% Información del documento
\title{Hoja 4: Ejercicio 18}
\subtitle{TEMA}
\author{Diego Rodríguez Cubero}
\institute{UCM}
\date{\today}

% Logo (opcional)
% \logo{\includegraphics[height=1cm]{images/logo.png}}

\begin{document}

% Diapositiva de título
\begin{frame}
    \titlepage
\end{frame}

% Tabla de contenidos
\begin{frame}{Contenidos}
    \tableofcontents
\end{frame}

% Sección 1
\section{Enunciado}
\begin{frame}{Enunciado del ejercicio}
    \begin{block}{Ejercicio 8}
        Dado un grafo no dirigido, conexo, valorado y en el que no hay dos aristas que tengan el    mismo coste, razonar con el mayor rigor posible que los algoritmos de Kruskal y Prim
        necesariamente calcularán el mismo árbol de recubrimiento de coste mínimo.
    \end{block}
\end{frame}

% Sección 2
\section{Idea de la solución}
\begin{frame}{Idea de la solución}
    \begin{itemize}
        \item Ambos algoritmos calculan un ARM, por lo que basta ver que hay un único ARM.
        \item Tendremos que comprobar que toda arista del ARM es la mínima que conecta dos componentes conexas.
    \end{itemize}
\end{frame}

% Sección 3
\section{Definiciones}
\begin{frame}{Definiciones}
    \begin{itemize}
        \item \textbf{Árbol de recubrimiento mínimo (ARM)}: Árbol que conecta todos los vértices del grafo con el coste total mínimo.
        \item En nuestro caso, un grafo $G$ no dirigido, conexo, valorado y sin aristas de igual coste.
    \end{itemize}
\end{frame}

\begin{frame}{Algoritmos de Kruskal y Prim}
    \begin{itemize}
        \item \textbf{Kruskal}: Añade aristas de menor coste que no formen ciclos.
        \item \textbf{Prim}: Añade aristas de menor coste que conecten un vértice del árbol con otro fuera de él.
    \end{itemize}
\end{frame}

% Sección 4
\section{Demostración}
\begin{frame}{Demostración por absurdo}
    \begin{enumerate}
        \item Supongamos que hay dos ARM distintos $A$ y $B$.
        \item Sea $e$ la arista de menor coste que pertenece a $A$ pero no a $B$ (tal arista existe ya que $A \nsubseteq B$, si no $B$ no podría ser un ARM al poder quitarle aristas hasta quedarse con $A$).
        \item Al eliminar $e$ de $A$, se forman dos componentes conexas $S$ y $\overline{S}$ (si no fuera así, $A$ tendría un ciclo).
        \item Como $B$ es conexo, existe una arista $f$ en $B$ que conecta $S$ y $\overline{S}$, y dicha arista no puede estar en $A$, ya que si hacemos el corte $A - \{e, f\}$, segun proposiciones vistas en clase, solo la menor arista que conecta $S$ y $\overline{S}$ puede pertenecer a un ARM, y esa sera solo una de las dos.
        \item Si ahora tomamos $A' = A - \{e\} + \{f\}$, obtenemos otro árbol de recubrimiento (ya que $f$ conecta $S$ y $\overline{S}$).
        
    \end{enumerate}
\end{frame}

\begin{frame}{Demostración por absurdo (cont.)}
    \begin{enumerate}
        \item Si el coste de $f$ fuera menor que el de $e$, entonces el coste de $A'$ sería menor que el de $A$, contradiciendo que $A$ es un ARM.
        \item Si el coste de $f$ fuera mayor que el de $e$, entonces el coste de $B' = B - \{f\} + \{e\}$ ($e \notin B$) sería menor que el de $B$, contradiciendo que $B$ es un ARM.
        \item Por tanto, el coste de $f$ debe ser igual al de $e$, lo cual es imposible ya que no hay dos aristas con el mismo coste, a no ser que $e = f$, lo cual es absurdo ya que $e \notin B$.
        \item Por tanto, nuestra suposición inicial es falsa y sólo existe un único ARM.
    \end{enumerate}
\end{frame}

% Sección 5
\section{Conclusión}
\begin{frame}{Conclusión}
    \begin{itemize}
        \item Dado que ambos algoritmos calculan un ARM y que sólo existe un único ARM, ambos algoritmos calcularán el mismo árbol de recubrimiento de coste mínimo, al ser la única opción posible.
    \end{itemize}
\end{frame}

\end{document}